\chapter{Policy Gradients and REINFORCE}

Chapters 5--11 mostly learn a policy \textbf{indirectly}. Policy iteration, SARSA, Q-learning, and DQN first estimate values $V$ or $Q$ (or $\hat{Q}$), then set $\pi$ by a greedy or $\varepsilon$-greedy rule on those values---for example $\pi(s)=\arg\max_a Q(s,a)$. That works when $\mathcal{A}$ is a small discrete set, as in the taxi's few turns: score each action, pick among them.

When $\mathcal{A}$ is not like that, the same pattern breaks down. Suppose the agent must set a continuous control---how hard to brake, how far to turn the wheel. A value-based learner would need $\hat{Q}(s,a)$ for every real-valued $a$, then search for the best one at every step; that is not the discrete $\arg\max$ of Chapters 8--11. Or suppose the setting is rock--paper--scissors: any fixed pure choice (always rock, always paper, or always scissors) is beaten as soon as the opponent notices; the agent needs a mixed strategy $\pi(a\mid s)$. $\varepsilon$-greedy (6.7) does randomize, but non-greedy actions share $\varepsilon/|\mathcal{A}|$ equally, so it cannot learn a specific mix such as $70\%$--$30\%$, and if two actions have nearly equal $Q$ and a small update swaps which of the two is greedy, almost all probability $1-\varepsilon$ suddenly moves from one to the other. In both cases what should be learned is the action-selection rule itself---not a policy reconstructed after each value update.

\textbf{Policy gradient} methods learn $\pi$ \textbf{directly}. The policy is written $\pi_\theta(a\mid s)$, a distribution over actions with parameters $\theta$, so a continuous action law or a specific mix is represented by $\theta$ itself, without $\arg\max_a\hat{Q}$. The policy is scored by a single number $J(\theta)$, and $\theta$ is improved by gradient ascent on $J$; the name says exactly that, the gradient of the objective with respect to the policy's parameters. The score is the expected episode return from the start state $s_0$:

\begin{equation}
J(\theta)
=
V^{\pi_\theta}(s_0)
=
\mathbb{E}_{\pi_\theta}\bigl[G_0\mid S_0=s_0\bigr].
\tag{9.1}
\end{equation}

The score is taken at the start state $s_0$ because $G_0$ already sums every reward in the episode, so $V^{\pi_\theta}(s_0)$ scores the whole episode. Raising it requires better actions in every state visited from $s_0$, since those actions determine $G_0$.

Gradient ascent needs $\nabla_\theta J$, but $J$ is an expectation over episodes whose distribution depends on the unknown environment, so $\nabla_\theta J$ cannot be computed exactly. \textbf{REINFORCE} (\citetxt{williams1992}{Williams, 1992}), the basic policy-gradient method, estimates it from sampled episodes: run $\pi_\theta$ for a full episode, compute the return that followed each action, and change $\theta$ to make each action taken more likely, weighted by that return. An action followed by a high return becomes more probable in the state where it was taken; it is reinforced, which gives the method its name. Like the Monte Carlo methods of Chapter 4, REINFORCE learns from complete episodes and needs neither a model of the environment nor a value function.

\section{Parameterized policies}

Raising $J(\theta)$ requires a concrete $\pi_\theta$. From Chapter 1, $\pi(a\mid s)=P(A_t=a\mid S_t=s)$. A \textbf{parameterized policy} makes that probability depend on a parameter vector $\theta$:

\[
\pi_\theta(a\mid s)
=
P(A_t=a\mid S_t=s;\,\theta).
\]

For each $s$, the numbers $\{\pi_\theta(a\mid s)\}_{a\in\mathcal{A}}$ must be nonnegative, sum to $1$, and change with $\theta$. That still needs a map from $(s,\theta)$ to a distribution over $\mathcal{A}$.

The softmax of Chapter 6 (6.9) is such a map: one real score per action becomes a distribution over $\mathcal{A}$. There the scores were action-values $Q(s,a)$. That coupling is the wrong object for policy gradients: $Q$ must track expected return and the Bellman backup, so its scale is tied to rewards and $\gamma$, and $\pi$ is only a byproduct of those values. What is needed instead is a parameterized state--action score whose only job is to set $\pi_\theta$---free to rise or fall as raising $J$ requires. Call that score a \textbf{preference} $h(s,a;\theta)$. Learning means updating $\theta$, which changes $h$, which changes $\pi_\theta$. Larger $h$ means larger probability after the softmax. Concretely, $\theta$ may be a table with one entry $\theta_{s,a}$ per pair ($h=\theta_{s,a}$), or the weights of a network that maps $s$ to $|\mathcal{A}|$ scores. For discrete $\mathcal{A}$, the parameterized policy is the softmax of these preferences:

\[
\pi_\theta(a\mid s)
=
\frac{\exp\bigl(h(s,a;\theta)\bigr)}
{\sum_{b\in\mathcal{A}}\exp\bigl(h(s,b;\theta)\bigr)}.
\]

For continuous $\mathcal{A}$ there is no finite list to softmax over, so $\pi_\theta(\cdot\mid s)$ is a continuous distribution. A common choice is Gaussian. Just as $\theta$ parameterizes preferences $h$ in the discrete case, here $\theta$ parameterizes a map from the state $s$ to the Gaussian's mean $\mu_\theta(s)$ and standard deviation $\sigma_\theta(s)>0$. The density of that policy is

\[
\pi_\theta(a\mid s)
=
\frac{1}{\sqrt{2\pi}\,\sigma_\theta(s)}
\exp\Biggl(
-\frac{\bigl(a-\mu_\theta(s)\bigr)^2}{2\sigma_\theta(s)^2}
\Biggr).
\]

Actions are sampled from that distribution.

Either form---softmax of preferences or Gaussian---gives a concrete $\pi_\theta$ the agent can sample from. Raising $J(\theta)$ is then gradient ascent on $\theta$:

\[
\theta
\leftarrow
\theta+\alpha\,\nabla_\theta J(\theta).
\]

The rest of this section derives a usable estimate of $\nabla_\theta J$ and turns it into the \textbf{REINFORCE} algorithm.

\section{Trajectory form of $J(\theta)$}

Raising $J$ by gradient ascent needs $\nabla_\theta J(\theta)$. By (9.1), $J(\theta)$ is the expected return of episodes that start at $s_0$ and are generated by $\pi_\theta$. A full episode is a trajectory \[
\tau
=
(S_0,A_0,R_1,S_1,A_1,R_2,\ldots,S_T).
\] Its probability under $\pi_\theta$ factors into the start-state distribution, the policy choices, and the environment transitions from Chapter 1: \begin{equation}
P(\tau;\theta)
=
P(S_0)
\prod_{t=0}^{T-1}
\pi_\theta(A_t\mid S_t)\,
P(S_{t+1},R_{t+1}\mid S_t,A_t).
\tag{9.2}
\end{equation} Only the $\pi_\theta$ factors depend on $\theta$. Write $R(\tau)$ for the return of $\tau$. That return is $G_0$ from Chapter 1, not the per-step reward $R_t$:

\[
R(\tau)
=
G_0
=
R_1+\gamma R_2+\gamma^2 R_3+\cdots.
\]

The expectation in (9.1) averages $R(\tau)$ over trajectories according to those probabilities. In the discrete case each trajectory has probability mass $P(\tau;\theta)$, so

\[
J(\theta)
=
\sum_{\tau} P(\tau;\theta)\,R(\tau).
\]

If trajectories form a continuous family, the sum becomes an integral with density $P(\tau;\theta)$:

\[
J(\theta)
=
\int P(\tau;\theta)\,R(\tau)\,d\tau.
\]

The integral is definite: it runs over all trajectories. The two forms are the same idea; the steps below use the integral notation, and the discrete case replaces integrals by sums.

\section{Policy gradient}

The return $R(\tau)$ is produced by the environment along the trajectory; it does not depend on the policy parameters $\theta$. Only the distribution $P(\tau;\theta)$ does. Differentiate the integral form of $J$ accordingly---$\nabla_\theta$ hits $P$, not $R$:

\[
\nabla_\theta J(\theta)
=
\int \bigl(\nabla_\theta P(\tau;\theta)\bigr)\,R(\tau)\,d\tau.
\]

For $u=P(\tau;\theta)>0$, $\nabla_\theta\log u=\frac{1}{u}\nabla_\theta u$, so $\nabla_\theta P(\tau;\theta)=P(\tau;\theta)\,\nabla_\theta\log P(\tau;\theta)$. Substitute:

\[
\nabla_\theta J(\theta)
=
\int P(\tau;\theta)\,\bigl(\nabla_\theta\log P(\tau;\theta)\bigr)\,R(\tau)\,d\tau.
\]

The integrand has the form $P(\tau;\theta)\,f(\tau)$, where \[
f(\tau)
=
\bigl(\nabla_\theta\log P(\tau;\theta)\bigr)\,R(\tau)
\] is a function of the trajectory. By definition of expectation under the density $P(\tau;\theta)$,

\begin{equation}
\nabla_\theta J(\theta)
=
\mathbb{E}_{\tau\sim\pi_\theta}\Bigl[
\bigl(\nabla_\theta\log P(\tau;\theta)\bigr)\,R(\tau)
\Bigr].
\tag{9.3}
\end{equation}

The subscript $\tau\sim\pi_\theta$ means the expectation is over trajectories generated by acting with $\pi_\theta$ (together with the environment): sample $\tau$ from $P(\tau;\theta)$. It remains to evaluate $\nabla_\theta\log P(\tau;\theta)$ from (9.2). The log of the product is a sum:

\begin{equation}
\begin{aligned}
\log P(\tau;\theta)
&=
\log P(S_0)
+
\sum_{t=0}^{T-1}\log\pi_\theta(A_t\mid S_t)
\\
&\qquad
+
\sum_{t=0}^{T-1}\log P(S_{t+1},R_{t+1}\mid S_t,A_t).
\end{aligned}
\tag{9.4}
\end{equation}

Differentiate with respect to $\theta$. The terms $P(S_0)$ and $P(s',r\mid s,a)$ do not depend on $\theta$, so

\begin{equation}
\nabla_\theta\log P(\tau;\theta)
=
\sum_{t=0}^{T-1}
\nabla_\theta\log\pi_\theta(A_t\mid S_t).
\tag{9.5}
\end{equation}

Substitute (9.5) into (9.3):

\begin{equation}
\nabla_\theta J(\theta)
=
\mathbb{E}_{\tau\sim\pi_\theta}\Biggl[
\Biggl(\sum_{t=0}^{T-1}\nabla_\theta\log\pi_\theta(A_t\mid S_t)\Biggr)
R(\tau)
\Biggr].
\tag{9.6}
\end{equation}

The gradient uses only $\pi_\theta$ and returns from sampled episodes.

\section{REINFORCE}

The factor $\nabla_\theta\log\pi_\theta(A_t\mid S_t)$ in (9.6) is a function of $(S_t,A_t)$ and $\theta$. Write \[
\psi_\theta(s,a)
=
\nabla_\theta\log\pi_\theta(a\mid s),
\] so (9.6) with $R(\tau)=G_0$ is \[
\nabla_\theta J(\theta)
=
\mathbb{E}_{\tau\sim\pi_\theta}\Biggl[
\sum_{t=0}^{T-1}
G_0\,\psi_\theta(S_t,A_t)
\Biggr].
\] The return $G_0$ multiplies every $\psi_\theta(S_t,A_t)$, and it includes rewards received before $A_t$ was taken. For each $t$, split $G_0$ into the rewards before time $t$ and the rest: \[
G_0
=
\sum_{k=0}^{t-1}\gamma^{k}R_{k+1}
+
\gamma^{t}G_t,
\qquad
G_t
=
\sum_{k=t}^{T-1}\gamma^{k-t}R_{k+1},
\] where $G_t$ is the return-to-go from time $t$. Substituting the split and applying linearity (A.4) separates the gradient into a past part and a future part: \[
\nabla_\theta J(\theta)
=
\underbrace{
\sum_{t=0}^{T-1}\sum_{k=0}^{t-1}\gamma^{k}\,
\mathbb{E}\bigl[R_{k+1}\,\psi_\theta(S_t,A_t)\bigr]
}_{\text{past}}
+
\underbrace{
\mathbb{E}_{\tau\sim\pi_\theta}\Biggl[
\sum_{t=0}^{T-1}
\gamma^{t}G_t\,\psi_\theta(S_t,A_t)
\Biggr]
}_{\text{future}}.
\] Every term of the past part is zero, by \eqref{eq:E.2} in \appref{E}. The reward $R_{k+1}$ with $k<t$ was received before $A_t$ was chosen, and the policy draws $A_t$ from $\pi_\theta(\cdot\mid S_t)$ whatever those earlier rewards were, so \[
\mathbb{E}\bigl[R_{k+1}\,\psi_\theta(S_t,A_t)\bigr]
=
0
\qquad\text{for } k<t.
\] Only the future part remains. Replacing $\psi_\theta(S_t,A_t)$ by its definition $\nabla_\theta\log\pi_\theta(A_t\mid S_t)$, \begin{equation}
\nabla_\theta J(\theta)
=
\mathbb{E}_{\tau\sim\pi_\theta}\Biggl[
\sum_{t=0}^{T-1}
\gamma^{t}G_t\,\nabla_\theta\log\pi_\theta(A_t\mid S_t)
\Biggr].
\tag{9.7}
\end{equation} In (9.7), each action is weighted only by the return that followed it, and every quantity can be read from a sampled episode. \textbf{REINFORCE} estimates this gradient from sampled episodes.

The same gradient can be written with action-values. By definition $Q^{\pi_\theta}(s,a)=\mathbb{E}[G_t\mid S_t=s,A_t=a]$, so $\mathbb{E}[G_t\mid S_t,A_t]=Q^{\pi_\theta}(S_t,A_t)$. Starting from (9.7), \[
\begin{aligned}
\nabla_\theta J(\theta)
&=
\mathbb{E}\Biggl[\sum_{t=0}^{T-1}\gamma^{t}G_t\,\psi_\theta(S_t,A_t)\Biggr]
\\
&=
\sum_{t=0}^{T-1}\gamma^{t}\,\mathbb{E}\bigl[G_t\,\psi_\theta(S_t,A_t)\bigr]
\quad\text{[A.4]}
\\
&=
\sum_{t=0}^{T-1}\gamma^{t}\,\mathbb{E}\Bigl[\mathbb{E}\bigl[G_t\,\psi_\theta(S_t,A_t)\bigm|S_t,A_t\bigr]\Bigr]
\quad\text{[A.8]}
\\
&=
\sum_{t=0}^{T-1}\gamma^{t}\,\mathbb{E}\Bigl[\psi_\theta(S_t,A_t)\,\mathbb{E}\bigl[G_t\bigm|S_t,A_t\bigr]\Bigr]
\quad\text{[A.5; $\psi_\theta(S_t,A_t)$ is fixed given $S_t,A_t$]}
\\
&=
\sum_{t=0}^{T-1}\gamma^{t}\,\mathbb{E}\bigl[\psi_\theta(S_t,A_t)\,Q^{\pi_\theta}(S_t,A_t)\bigr]
\quad\text{[definition of $Q^{\pi_\theta}$]}
\\
&=
\mathbb{E}\Biggl[\sum_{t=0}^{T-1}\gamma^{t}Q^{\pi_\theta}(S_t,A_t)\,\psi_\theta(S_t,A_t)\Biggr].
\quad\text{[A.4]}
\end{aligned}
\] Replacing $\psi_\theta(S_t,A_t)$ by its definition $\nabla_\theta\log\pi_\theta(A_t\mid S_t)$, \begin{equation}
\nabla_\theta J(\theta)
=
\mathbb{E}_{\tau\sim\pi_\theta}\Biggl[
\sum_{t=0}^{T-1}
\gamma^{t}Q^{\pi_\theta}(S_t,A_t)\,\nabla_\theta\log\pi_\theta(A_t\mid S_t)
\Biggr].
\tag{9.7a}
\end{equation} Form (9.7a) needs $Q^{\pi_\theta}$, which is unknown. In (9.7), $G_t$ is a single sample whose conditional mean is $Q^{\pi_\theta}(S_t,A_t)$.

The same gradient can also be written without an explicit sum over time: average $Q^{\pi_\theta}(s,a)\,\nabla_\theta\log\pi_\theta(a\mid s)$ under the discounted state--action visits of $\pi_\theta$ from $s_0$. That identity is the \textbf{policy gradient theorem} (\citetxt{sutton2000}{Sutton et al., 2000}); \appref{E} derives it from the Bellman equation for $V^{\pi_\theta}$ and shows that it equals (9.7a).

The expectation in (9.7) cannot be computed, but it can be approximated by an average over episodes $\tau^{(1)},\ldots,\tau^{(N)}$ generated by $\pi_\theta$, as in the Monte Carlo methods of Chapter 4: \[
\nabla_\theta J(\theta)
\approx
\frac{1}{N}\sum_{i=1}^{N}
\sum_{t=0}^{T_i-1}
\gamma^{t}G_t^{(i)}\,\nabla_\theta\log\pi_\theta\bigl(A_t^{(i)}\mid S_t^{(i)}\bigr),
\] where the superscript $(i)$ marks quantities from episode $i$. REINFORCE updates $\theta$ after every single episode, so $N=1$: the average over episodes reduces to the one episode's sum over $t$, and the superscript is dropped. Generate that episode, compute each $G_t$, and put the estimate in place of $\nabla_\theta J$ in the gradient-ascent step $\theta\leftarrow\theta+\alpha\nabla_\theta J(\theta)$: \begin{equation}
\theta
\leftarrow
\theta
+
\alpha
\sum_{t=0}^{T-1}
\gamma^{t}G_t\,\nabla_\theta\log\pi_\theta(A_t\mid S_t).
\tag{9.8}
\end{equation}

Comparing (9.8) with the gradient-descent step $\theta\leftarrow\theta-\alpha\nabla_\theta L(\theta)$ gives the loss, since each $G_t$ is a fixed number computed from the episode: \begin{equation}
L(\theta)
=
-\sum_{t=0}^{T-1}
\gamma^{t}G_t\,\log\pi_\theta(A_t\mid S_t).
\tag{9.9}
\end{equation}

\begin{implnote}[Algorithm (REINFORCE)]
Input: differentiable policy $\pi_\theta(a\mid s)$; step size $\alpha$; discount $\gamma$; number of episodes $M$.

Output: trained policy $\pi_\theta$.

\begin{enumerate}
\item Initialize $\theta$ (e.g. randomly).
\item Loop for episode $=1,2,\ldots,M$:
\begin{itemize}
\item Generate an episode $\tau=(S_0,A_0,R_1,\ldots,S_T)$ by sampling $A_t\sim\pi_\theta(\cdot\mid S_t)$ until termination.
\item For each $t=0,1,\ldots,T-1$:
\begin{itemize}
\item $G_t\leftarrow\sum_{k=t}^{T-1}\gamma^{k-t}R_{k+1}$.
\end{itemize}
\item $\theta\leftarrow\theta+\alpha\sum_{t=0}^{T-1}\gamma^{t}G_t\,\nabla_\theta\log\pi_\theta(A_t\mid S_t)$.
\end{itemize}
\item Return $\pi_\theta$.
\end{enumerate}

\end{implnote}

\lstinputlisting{../code/ch09/reinforce.py}


REINFORCE is on-policy and Monte Carlo: each update uses full returns from episodes generated by the current $\pi_\theta$. The estimator (9.7) is unbiased but can have high variance because $G_t$ sums many noisy rewards. Chapter 10 improves that weight the way Chapters 6--7 improved value targets: first center $G_t$ with a baseline, then replace it by a TD error.
